% ===================================================================== PORTADA
\thispagestyle{empty}
\begin{tikzpicture}[remember picture,overlay]
  \fill[tinta] (current page.south west) rectangle (current page.north east);
  \fill[colY] (current page.north west) rectangle ([xshift=0.9cm]current page.south west);

  % --- texto
  \node[anchor=north west,font=\sffamily\bfseries\footnotesize,text=marca]
    at ([xshift=2.6cm,yshift=-2.6cm]current page.north west)
    {PROBABILIDAD \quad\textbullet\quad FACULTAD DE INGENIER\'IA \quad\textbullet\quad SEMESTRE 2027-1};
  \node[anchor=north west,font=\sffamily\bfseries\fontsize{58}{58}\selectfont,text=colY]
    at ([xshift=2.5cm,yshift=-3.6cm]current page.north west) {TEMA 3};
  \node[anchor=north west,font=\sffamily\bfseries\fontsize{29}{34}\selectfont,text=white,text width=17cm,align=left]
    at ([xshift=2.6cm,yshift=-6.1cm]current page.north west) {Variables aleatorias conjuntas};
  \node[anchor=north west,font=\sffamily\fontsize{15}{19}\selectfont,text=white!80!tinta,text width=13cm,align=left]
    at ([xshift=2.6cm,yshift=-7.7cm]current page.north west)
    {Gu\'ia paso a paso para el examen:\\ de cero al 10, con todos los moldes de la serie.};

  % --- gr\'afico: una distribuci\'on conjunta con sus dos marginales
  \begin{scope}[shift={([xshift=5.2cm,yshift=4.6cm]current page.south west)}]
    \foreach \i in {0,...,5}{
      \gdef\acc{0}
      \foreach \j in {0,...,5}{
        \pgfmathsetmacro\v{exp(-((\i-2.4)^2+(\j-2.7)^2-1.05*(\i-2.4)*(\j-2.7))/2.7)}
        \pgfmathparse{\acc+\v}\xdef\acc{\pgfmathresult}
        \pgfmathtruncatemacro\p{100*\v}
        \fill[marca!\p!tinta2] (\i*1.5,\j*1.5) rectangle ++(1.38,1.38);
      }
      \fill[colX] (\i*1.5,9.25) rectangle ++(1.38,\acc*0.62);
    }
    \foreach \j in {0,...,5}{
      \gdef\acc{0}
      \foreach \i in {0,...,5}{
        \pgfmathsetmacro\v{exp(-((\i-2.4)^2+(\j-2.7)^2-1.05*(\i-2.4)*(\j-2.7))/2.7)}
        \pgfmathparse{\acc+\v}\xdef\acc{\pgfmathresult}
      }
      \fill[colY] (9.25,\j*1.5) rectangle ++(\acc*0.62,1.38);
    }
    \node[anchor=west,font=\sffamily\bfseries\small,text=colX!55!white] at (0,11.75) {marginal de $X$: sumo cada columna};
    \node[anchor=center,font=\sffamily\bfseries\small,text=colY!60!white,rotate=-90] at (12.1,4.44) {marginal de $Y$: sumo cada rengl\'on};
    \node[anchor=north,font=\sffamily\bfseries\small,text=marca] at (4.44,-0.25) {la conjunta $f_{XY}(x,y)$: una probabilidad para cada pareja};
  \end{scope}

  \node[anchor=south west,font=\sffamily\small,text=white!70!tinta,text width=15cm,align=left]
    at ([xshift=2.6cm,yshift=1.5cm]current page.south west)
    {Hecha con tus apuntes, tu formulario y la Serie 3 \quad\textbullet\quad
     recetas, ejemplos resueltos paso a paso y un simulacro con soluci\'on};
\end{tikzpicture}
\clearpage

% ===================================================================== \'INDICE
\thispagestyle{fancy}
\markboth{Contenido}{}
{\sffamily\bfseries\fontsize{24}{28}\selectfont\color{tinta}Contenido}\par
\vspace{-4pt}{\color{colY}\rule{2.2cm}{3pt}}\par\medskip
\makeatletter\@starttoc{toc}\makeatother
